\chapter{Introduction}
\label{ch:introduction}

% Une introduction de mémoire présente généralement le contexte, la question
% précise, la contribution et un guide de lecture. Pour un rapport long,
% quatre à six pages suffisent souvent ; il s'agit d'un conseil de
% planification et non d'une règle LMFI.

Les graphes orientés modélisent des relations à sens unique : une dépendance peut
relier une tâche à une autre, ou une transition peut conduire d'un état de
programme à son successeur. Une première question consiste à déterminer si un
sommet est accessible depuis un autre en suivant zéro ou plusieurs arêtes. Ce
problème élémentaire relie ici définitions mathématiques, démonstrations,
algorithmes, figures, tableaux et références dans un même exemple. Des
introductions classiques figurent notamment dans \cite{diestel2017,cormen2022}.

\section{Question et contribution}

L'exemple répond à la question suivante : comment définir l'accessibilité dans
un graphe orienté fini, démontrer sa transitivité et la calculer à partir d'une
matrice d'adjacence ? La contribution reste volontairement modeste. Nous
formalisons la relation au \cref{ch:background}, la calculons au
\cref{ch:main-work} et étudions un cas concret au \cref{ch:evaluation}.

\section{Organisation de l'exemple}

La séparation entre préliminaires, travail principal et évaluation illustre une
structure de chapitres maintenable. Un mémoire théorique peut diviser le travail
principal en plusieurs chapitres de résultats ; un mémoire appliqué peut diviser
l'évaluation en protocole, résultats et discussion. Les fichiers sources rendent
ces deux adaptations locales et explicites.
